\section{Introduction}

Federated learning coordinates multiple clients without centralizing their private training data. Model-heterogeneous federated learning (HFL) removes the additional assumption that clients share an architecture, so collaboration must proceed through transferable objects such as predictions, prototypes, or public-data responses. Representative approaches include public-logit distillation in FedDF \citep{lin2020feddf}, prototype exchange in FedProto \citep{tan2022fedproto}, robust collaboration under noisy labels in RHFL \citep{fang2022rhfl}, and corruption-aware heterogeneous learning in AugHFL and RAHFL \citep{fang2023aughfl,fang2025rahfl}. In the latter line, corruption is primarily a nuisance that should not change the semantic label.

Corruption can play a more damaging role when it is distributed unevenly across classes. Consider a client that observes one class mainly with blur and another mainly with noise, while a second client exhibits a different association. The operator is then predictive of the task label within each client. A model can exploit this association even if its predictions are locally consistent under stochastic augmentation. We call this controlled failure mode \textbf{class--corruption entanglement in HFL (CLE-HFL)}. It is not equivalent to low corrupted accuracy: image degradation, label skew, capacity, and non-IID optimization can all lower accuracy without causing predictions to follow a class--corruption binding.

This distinction calls for a paired diagnostic. For a fixed semantic source, we intervene on the corruption operator and measure whether predictive probability moves toward the classes associated with that operator during training. The resulting statistic, \textbf{Directional Shortcut Alignment (DSA)}, is evaluated against a binding fixed before model inspection. Source pairing cancels operator-invariant source terms, while a shuffled-binding null tests whether the response follows the registered association rather than an arbitrary operator preference. DSA does not recover the model's complete internal causal mechanism; under explicit semantic-preservation and evaluation-isolation assumptions, it estimates a binding-direction operator-response contrast. We further show why neither within-operator consistency nor an unconstrained pseudo-environment proxy can certify that this contrast has disappeared.

The same design allows us to ask where CLE emerges. Across three matched training seeds, the DSA contrast between strong CLE and its control is already large without communication, whereas the additional HFL communication contrast is small on average and changes sign across seeds. We refer to this empirical pattern as \textbf{predominantly local-first}. This finding does not make the problem non-federated: the shortcut is induced by client-specific federated data distributions, appears in heterogeneous models that later exchange knowledge, and can affect the information presented to the collaboration mechanism. It does, however, suggest that the first intervention point should be the client objective rather than a new server protocol.

The local-first result points to a client-side intervention with two components. A coarse, taxonomy-assisted \textbf{Public Environment Witness (PEW)} is trained on public images with synthetic corruption-family labels and then frozen. On private fit data, \textbf{Balanced Environment Risk (BER)} uses the resulting pseudo-environments to compress support-induced risk-mass differences within each task class. PEW supplies side information; BER is the optimization rule. Neither component reads private true corruption metadata, and the communication protocol remains unchanged. This construction is compared with ERM, tail-risk optimization, and a shared-PEW GroupDRO control so that access to pseudo-groups is not confused with the effect of BER itself.

Our study separates claims that are often conflated. Fixed-communication comparisons test the local objective; alternative binding maps test dependence on a particular association; a second communication base tests whether DSA suppression is tied to one collaboration rule; and a mechanism-family table asks whether standard HFL or corruption-robust methods remove CLE on their own. DSA and task utility are reported separately because reducing a shortcut need not improve every accuracy metric. The contributions are:

\begin{enumerate}
\item \textbf{A controlled HFL failure mode.} We formulate client-specific class--corruption bindings and separate binding-directed shortcut use from generic corruption difficulty.
\item \textbf{A paired diagnostic with an explicit identification boundary.} We introduce source-paired operator evaluation and DSA, establish its cancellation and binding-specificity properties, use source-level inference, and prove that unconstrained proxy statistics do not identify true CLE dependence.
\item \textbf{A local-versus-communication attribution study.} A matched factorial reveals a predominantly local-first pattern in the evaluated setting and quantifies the smaller communication add-on without interpreting the ratio as sample-level mediation.
\item \textbf{A targeted mitigation and its empirical boundaries.} PEW+BER adjusts class-conditional pseudo-environment support without redesigning communication. We test it against matched optimization controls and across controlled bindings and communication bases, while treating taxonomy dependence and non-uniform utility as limitations rather than hiding them.
\end{enumerate}

\ifdraftskeleton
\evidenceneeded{remaining moderate Axis-II rows, second private dataset, and bounded taxonomy stress.}
\fi
